%
% Der euklidische Algorithmus für ganz Zahlen
%
\subsection{Der euklidische Algorithmus für ganze Zahlen
\label{buch:ratint:euklid:subsection:ganzezahlen}}
In diesem Abschnitt wird der euklidische Algorithmus für ganze
Zahlen entwickelt.
Diese Form des Algorithmus ist seit dem Altertum bekannt und wird
oft bereits in der Mittelschule ein sowohl historisch
wie auch mathematisch bedeutsames Beispiel eines Algorithmus
unterrichtet.
Im Abschnitt~\ref{buch:ratint:euklid:subsection:polynome} wird
der Algorithmus auf Polynome verallgemeinert.

%
% Der Grundalgorithmus
%
\subsubsection{Der Grundalgorithmus}
Seien $a>b$ natürliche Zahlen und $q$ und $r$ so gewählte
natürliche Zahlen, dass $a=bq+r$ mit $r<b$.
Der grösste gemeinsame Teiler $g=\operatorname{ggT}(a,b)$ teilt dann
sowohl $a$ als auch $b$, er muss daher auch $r$ teilen.
Wiederholt man die Rechnung mit $b$ anstelle von $a$ und $r$ anstelle
von $b$ folgt, dass $g$ erneut den Rest teilt.
Da der Rest immer kleiner wird, gibt es einen Moment, bei dem
er erstmals verschwindet.
Der Divisor $b$ ist dann der grösste gemeinsame Teiler.

\begin{beispiel}
\label{buch:ratint:euklid:bsp:euklidzahlen}
Für die Zahlen $a=36$ und $b=28$ ergibt sich die Rechnung
\begin{center}
\input{chapters/060-ratint/fig/tabular-euklid.tex}
\end{center}
aus der man $\operatorname{ggT}(36,28)=4$ ablesen kann.
\end{beispiel}

%
% Der erweiterte Algorithmus
%
\subsubsection{Der erweiterte Algorithmus}
Wir analysieren den euklidischen Algorithmus nochmals etwas im Detail.
Dazu führen wir die folgende Notation ein.
Wir beginnen mit $a_0=a$ und $b_0=b$.
\index{euklidischer Algorithmus!erweitert}%
Wir konstruieren dann $q_0$ und $r_0$ mit $a_0=q_0b_0+r_0$ und 
setzen anschliessend $a_1 = b_0$ und $b_1 = r_0$.
Dieser Prozess wird für $k>0$ gemäss
\begin{align*}
a_k &= b_{k-1} &
b_k &= r_{k-1} &
&\text{und}&
a_k &= q_k b_k + r_k 
&&\Rightarrow&
r_k &= a_k - q_kb_k
\end{align*}
wiederholt, bis $r_{n+1}=0$ ist.
Der grösste gemeinsame Teiler ist dann $r_n$.
Rechnet man ausgehend von $r_n$ rückwärts, findet man
\begin{align*}
r_n &= a_n - q_n b_n
\\
    &= b_{n-1} - q_n r_{n-1}
\\
    &= b_{n-1} - q_n ( a_{n-1} - q_{n-1} b_{n-1} )
     = q_n a_{n-1} + (1 + q_n q_{n-1}) b_{n-1}
\\[-4pt]
    &\phantom{1}\vdots
\intertext{Durch wiederholte Anwendung der Definitionsformeln
für $a_k$ und $b_k$ finden wir schliesslich ganze Zahlen $s$ und $t$
mit der Eigenschaft}
    &= s a_0 + t b_0.
\end{align*}

\begin{beispiel}
Für das Beispiel~\ref{buch:ratint:euklid:bsp:euklidzahlen} folgt
aus der zweiten Zeile der Tabelle
\[
\operatorname{ggT}(a,b)
=
4
= 
28 - 3\cdot 8
\]
Da 8 der Rest der vorangegangenen Zeile ist, kann man ihn durch
$36-1\cdot 28$ ersetzen und erhält
\[
4
=
28 - 3 \cdot (36-28)
=
4\cdot 28 - 3\cdot 36.
\]
Damit sind die Zahlen $s=-3$ und $t=4$ gefunden.
\end{beispiel}

Der erweiterte euklidische Algorithmus kann also nicht nur den 
grössten gemeinsamen Teiler von $a$ und $b$ finden, sondern zusätzlich
Zahlen $s$ und $t$ derart, dass
\[
sa + tb
=
\operatorname{ggT}(a,b)
\]
gilt.
Die Berechnung nach Abschluss ist jedoch eher umständlich.
Gewünscht ist eine Erweiterung des Grundalgorithmus, die die Zahlen
$s$ und $t$ gleich mitberechnet.

\begin{satz}
Gegeben sind die natürlichen Zahlen $a$ und $b$.
Der folgende Algorithmus berechnet den grössten gemeinsamen Teiler $g$
und Zahlen $s$ und $t$ mit $g=sa+tb$.
\begin{enumerate}
\item
Setze $a_0=a$, $b_0=b$, $k=0$ und $s_{-2}=1$, $t_{-2}=0$, $s_{-1}=0$
und $t_{-1}=1$
\item
Finde $q_k$ und $q_k$ mit $a_k= q_kb_k+r_k$.
\item
Berechne
$s_k=s_{k-2}-q_ks_{k-1}$ 
und
$t_k=t_{k-2}-q_kt_{k-1}$.
\end{enumerate}
Für alle $k$ gilt
$r_k = s_ka+t_kb$
Für den grössten Index $n$, für den $r_n\ne 0$ ist, folgt
\[
\operatorname{ggT}(a,b)
=
s_na+t_nb
\qquad\text{und}\qquad
0
=
s_{n+1}a+t_{n+1}b.
\]
\end{satz}

\begin{proof}
Wir schreiben die Division als Matrizenprodukt:
\[
\begin{pmatrix}
b_k\\
r_k
\end{pmatrix}
=
\begin{pmatrix*}[r]
b_k\\
a_k-q_kb_k
\end{pmatrix*}
=
\underbrace{
\begin{pmatrix}
 0 & 1   \\
 1 &-q_k
\end{pmatrix}
}_{\displaystyle = Q(q_k)}
\begin{pmatrix}
a_k\\
b_k
\end{pmatrix}.
\]
Wiederholte Anwendung ergibt
\[
\begin{pmatrix}
r_n\\
0
\end{pmatrix}
=
Q(q_{n+1})
Q(q_n)
\cdots
Q(q_1)
Q(q_0)
\begin{pmatrix}
a\\ b
\end{pmatrix}.
\]
Das Produkt der Matrizen enthält in der ersten Zeile die gesuchten
Zahlen $s_n$ und $t_n$ und in der zweiten Zeile die Zahlen
$s_{n+1}$ und $t_{n+1}$.

Da das Matrizenprodukt assoziativ ist, kann das Produkt von rechts
her aufgelöst werden.
Dazu definieren wir die Matrix
\[
S_k
= 
\begin{pmatrix}
s_{k-1} & t_{k-1} \\
s_{k}   & t_{k} 
\end{pmatrix}
\]
mit der initialen Matrix
\[
S_{-1}
=
\begin{pmatrix}
s_{-2} & t_{-2} \\
s_{-1} & t_{-1} 
\end{pmatrix}
=
\begin{pmatrix}
1 & 0\\
0 & 1
\end{pmatrix}.
\]
Die Multiplikation mit $Q(q_k)$ mit $S_{k-1}$ erzeugt $S_k$, es gilt
daher die Rekursionsformel
\begin{align*}
S_k
&= 
\begin{pmatrix}
s_{k-1} & t_{k-1} \\
s_{k}   & t_{k} 
\end{pmatrix}
=
Q(q_k)
S_{k-1}
=
Q(q_k)
\begin{pmatrix}
s_{k-2} & t_{k-2} \\
s_{k-1} & t_{k-1} 
\end{pmatrix}
\\
&=
\begin{pmatrix}
0&1\\
1&-q_k
\end{pmatrix}
\begin{pmatrix}
s_{k-2} & t_{k-2} \\
s_{k-1} & t_{k-1} 
\end{pmatrix}
\\
&=
\begin{pmatrix*}[r]
s_{k-1}            & t_{k-1}           \\
s_{k-2}-q_ks_{k-1} & t_{k-2}-q_kt_{k-1} 
\end{pmatrix*}.
\end{align*}
In der zweiten Zeile stehen die Rekursionsformeln für die
Zahlen $s_k$ und $t_k$.
\end{proof}

\begin{beispiel}
\label{buch:ratint:euklid:bsp:erweitert}
Wir führen den erweiterten euklidischen Algorithmus für $a=43$ und $b=47$
durch:
\begin{center}
\begin{tabular}{|>{$}r<{$}|>{$}r<{$}>{$}r<{$}|>{$}r<{$}>{$}r<{$}|>{$}r<{$}>{$}r<{$}|}
\hline
  k &  a_k &  b_k &  q_k &  r_k &  s_k &  t_k \\
\hline
 -2 &      &      &      &      &    1 &    0 \\
 -1 &      &      &      &      &    0 &    1 \\
  0 &   43 &   47 &    0 &   43 &    1 &    0 \\
  1 &   47 &   43 &    1 &    4 &   -1 &    1 \\
  2 &   43 &    4 &   10 &    3 &   11 &  -10 \\
  3 &    4 &    3 &    1 &
\smash{\rlap{\raisebox{-14pt}{\begin{tikzpicture}[>=latex,thick]
\fill[color=darkred!20] (0,0) rectangle (2.5,0.8);
\end{tikzpicture}}}}%
                   \phantom{0}1 &  -12 &   11 \\
  4 &    3 &    1 &    3 &    0 &   47 &  -43 \\
\hline
\end{tabular}
\end{center}
In der letzten Zeile wird der Rest $0$, der grösste gemeinsame Teiler ist
also $\operatorname{ggT}(a,b) = r_3=1$.
In der Zeile $k=3$ kann man die beiden Zahlen $s=s_3=-12$ und $t=t_3=11$
ablesen, und tatsächlich ist
\[
-12\cdot 43 + 11\cdot 47 =  1
\qquad\text{und}\qquad
47\cdot 43 - 43\cdot 47 = 0.
\qedhere
\]
\end{beispiel}

Das Beispiel illustriert auch, dass die im Grundalgorithmus formulierte
Bedingung, dass $a>b$ sein muss, unwichtig ist.
Wenn nämlich $a<b$ ist, dann ist der erste Quotient $q_0=0$ und der erste
Rest $r_0=a$.
Daraus ergibt sich dann $a_1=b$ und $b_1=a$, d.~h.~die beiden Zahlen $a$
und $b$ haben die Plätze getauscht.
Auch die Zahlen $s_0=1$ und $t_0=0$ tauschen die Plätze der Anfangsbedingungen
für die Folgen $s_k$ und $t_k$.

%
% Anwendung des erweiterten Algorithmus
%
\subsubsection{Anwendung des erweiterten Algorithmus}
Der erweiterte euklidische Algorithmus kann dazu verwendet werden,
die multiplikative inverse in einem endlichen Körper zu finden.
Sei $p$ eine Primzahl und $\mathbb{F}_p$ die Menge der Reste modulo $p$.
Die multiplikative Inverse eines Elementes $a\in\mathbb{F}_p$ kann
jetzt wie folgt gefunden werden.
Da $p$ ein Primzahl ist, sind $a$ und $p$ teilerfremd.
Der erweiterte euklidische Algorithmus liefert zwei Zahlen $s$ und $t$
mit
\[
sa + tp = \operatorname{ggT}(a,p) = 1.
\]
Modulo $p$ wird diese Gleichung zu
\[
sa \equiv 1 \mod p.
\]
Der Rest von $s$ modulo $p$ ist daher die multiplikative Inverse $a$.

\begin{beispiel}
Für $p=47$ und $a=43$ ergibt der in
Beispiel~\ref{buch:ratint:euklid:bsp:erweitert}
durchgeführte, erweiterte euklidische Algorithmus
\[
-12 \cdot 43
+
11\cdot 47
= 1.
\]
Die multiplikative Inverse von $43$ im Körper $\mathbb{F}_{47}$ ist
daher $-12\equiv 35\mod 47$.
\index{multiplikative Inverse in F_p@multiplikative Inverse in $\mathbb{F}_p$}%
Tatsächlich ist 
\[
35\cdot 43 = 1505 = 32\cdot 47 + 1,
\]
wie behauptet.
\end{beispiel}

